\lysh{14604}{3}

\begin{alist}
\item Το πεδίο ορισμού της συνάρτησης είναι σε μορφή διαστήματος το εξής: $A = (-\infty, 10)$.

\item Έχουμε:
\begin{gather*}
f(-1) = 2(-1) - 5 = -7, \\
f(3) = 2 \cdot 3 - 5 = 1\text{και} \\
f(5) = 5^2 = 25.
\end{gather*}

\item Για να διέρχεται η γραφική παράσταση της $f$ από την αρχή των αξόνων, πρέπει $f(0) = 0$.
Όμως $f(0) = 2 \cdot 0 - 5 = -5 \ne 0$, οπότε η γραφική παράσταση της $f$ δεν διέρχεται από την αρχή
των αξόνων.

\item Έχουμε δυο περιπτώσεις:
\begin{bulletlist}
\item $y = f(x) = 21 \Leftrightarrow 2x - 5 = 21 \Leftrightarrow 2x = 26 \Leftrightarrow x = 13 \notin (-\infty, 3]$
\item $y = f(x) = 21 \Leftrightarrow x^2 = 21 \Leftrightarrow x = \pm \sqrt{21}$. Από τις δυο αυτές τιμές του $x$ δεκτή είναι η $\sqrt{21}$ 
διότι $\sqrt{21} \in (3, 10)$.
\end{bulletlist}
Κατά συνέπεια το ζητούμενο σημείο είναι $(\sqrt{21}, 21)$.
\end{alist}
